\documentclass[11pt]{article}
\usepackage[utf8]{inputenc}
\usepackage[T1]{fontenc}
\usepackage{booktabs}
\usepackage{amsmath}
\usepackage{geometry}
\geometry{margin=2.5cm}

\title{Causal Fairness Analysis Report}
\author{Generated by FairMind and an LLM}
\date{2026-08-22}

\begin{document}
\maketitle

\section{Analysis setup}

\begin{tabular}{ll}
\toprule
Dataset & student\_mat \\
Protected attribute ($X$) & S2\_sex \\
Comparison & F $\to$ M \\
Outcome ($Y$) & T\_grade (1) \\
Mediator ($W$) & studytime \\
Confounder ($Z$) & Medu \\
\bottomrule
\end{tabular}

\section{Causal effects (FairMind, exact values)}

The five effects below are computed by FairMind through exact inference on the
fitted Bayesian network. These values are final and must not be recomputed or
altered.

\begin{tabular}{lr}
\toprule
Effect & Value \\
\midrule
Total Variation (TV) & 0.057084 \\
Total Effect (TE) & 0.049199 \\
Spurious Effect (SE) & 0.007886 \\
Direct Effect (DE) & 0.070005 \\
Indirect Effect (IE, decomposition form: TE = DE + IE) & -0.020806 \\
\bottomrule
\end{tabular}

\section{Qualitative interpretation}

\subsection{Total effect and spurious component (TV, TE, SE)}

The total disparity between the two groups is 0.057084, indicating a moderate difference in outcomes. After removing the confounding effect of the mediator, the genuine causal effect (TE) is 0.049199, which is slightly smaller than the total variation. The spurious effect (SE) of 0.007886, driven by the confounder Z, is relatively small but still contributes to the observed difference.

\subsection{Direct effect (DE)}

The direct effect (DE) of 0.070005 suggests a significant direct discrimination against the protected group, as it indicates a direct impact of the protected attribute on the outcome, bypassing the mediator. This direct effect is notably larger than the total effect, highlighting the substantial influence of direct discrimination.

\subsection{Indirect effect (IE)}

The indirect effect (IE) of -0.020806 shows that the mediator channel offsets the direct effect, reducing the overall impact of the direct discrimination. This negative IE means that the indirect pathway through the mediator mitigates the direct effect, resulting in a total effect that is smaller than the direct effect alone.

\section{Recap Questions}

\begin{enumerate}
\item Is there direct discrimination against the protected group? \textbf{Answer:} YES
\item Does the indirect effect through the mediator mitigate the total effect? \textbf{Answer:} YES
\item Does the observed total effect exceed the practical relevance threshold? \textbf{Answer:} YES
\item Does the spurious component (SE) contribute substantially to the observed difference? \textbf{Answer:} NO
\item Is the direct effect the dominant component compared to the indirect effect? \textbf{Answer:} YES
\end{enumerate}

\vspace{1em}
\noindent\footnotesize
The recap answers follow deterministic threshold rules applied to the values above: direct discrimination when $|\mathrm{DE}| > 0.05$; a mediator channel counted as negligible when $|\mathrm{IE}| < 0.005$; practical relevance when $|\mathrm{TV}| > 0.05$; a substantial spurious component when $|\mathrm{SE}| / |\mathrm{TV}| > 0.1$.

\end{document}